\documentclass[../../../include/open-logic-section]{subfiles}

\begin{document}

\olfileid{sth}{replacement}{ref}
\olsection{جایگزینی و بازتاب}

آخرین کوشش ما برای توجیه جایگزینی از راه \stagesinex، با نتیجه‌ای ژرف و زیبا آغاز می‌شود:\footnote{یادآوری می‌کنیم که همهٔ فرمول‌ها می‌توانند پارامتر داشته باشند، مگر آن‌که خلافش به‌صراحت گفته شود.}
%در اینجا از خط بالایی، مانند $x_1, \ldots, x_n$، استفاده می‌کنیم تا
%دنبالهٔ متغیرها را به‌اختصار نمایش دهیم.
\begin{thm}[طرحوارهٔ بازتاب]\ollabel{reflectionschema}
برای هر فرمول $\phi$:
\[
\forall \alpha \exists \beta > \alpha (\forall x_1 \ldots, x_n \in
V_\beta)(\phi(x_1, \ldots, x_n) \liff \phi^{V_\beta}(x_1, \ldots, x_n))
\]
\end{thm}
\noindent
همان‌گونه که در \olref[sth][replacement][strength]{formularelativization} دیدیم، $\phi^{V_\beta}$ از محدودکردن هر سورِ
$\phi$ به مجموعهٔ~$V_\beta$ به دست می‌آید. بنابراین، به‌طور شهودی،
بازتاب می‌گوید: اگر $\phi$ در تمام سلسله‌مراتب صادق باشد، آن‌گاه
$\phi$ در شمار دلخواهی از \emph{پاره‌های آغازین} سلسله‌مراتب نیز
صادق است.

\citet{Montague1961} و \citet{Levy1960} نشان دادند که (صورت‌بندی‌های
مناسبِ) جایگزینی و بازتاب، به پیمانهٔ~$\Z$، هم‌ارزند؛ پس افزودن هر
یک از آن‌ها $\ZF$ را به دست می‌دهد. (این نتایج را در
\olref[sth][replacement][refproofs]{sec} اثبات می‌کنیم.) با توجه به
این هم‌ارزی، شاید بتوان امید داشت که بازتاب و جایگزینی را با اتکا به
\stagesinex{} چنین توجیه کنیم: بنا بر \stagesinex، سلسله‌مراتب باید
بسیار بسیار بلند باشد؛ در واقع چنان بلند که هیچ سخنی که بتوانیم
درباره‌اش بگوییم برای کران‌دارکردن ارتفاع آن کافی نباشد. این را
می‌توان چنین فهمید که اگر جمله‌ای مانند~$\phi$ در تمام سلسله‌مراتب
صادق باشد، در شمار دلخواهی از پاره‌های آغازین آن نیز صادق است.
و این دقیقاً همان بازتاب است.

باز هم به نظر می‌رسد این کوششی واقعاً نویدبخش برای ارائهٔ توجیهی
درونی از جایگزینی باشد. اما مطالبی که باید درباره‌اش گفت بسیار بیش
از گنجایش اینجاست. اکنون خودتان باید داوری کنید که آیا این کوشش
موفق است یا نه.\footnote{البته شاید بخواهید مطالعه را با
\citet[95--100]{Incurvati2020} ادامه دهید.}

\end{document}
